\documentclass{ltjsarticle}
\usepackage{ascmac}
\usepackage{boites}
\begin{document}
\section{\LaTeX の環境を使ってみよう。}
ちょっと長い文章を枠の中に書いてみましょう。
\begin{breakbox}
「天は人の上に人を造らず人の下に人を造らず」と言えり。されば天より人を生ずるには、万人は万人みな同じ位にして、生まれながら貴賤きせん上下の差別なく、万物の霊たる身と心との働きをもって天地の間にあるよろずの物を資とり、もって衣食住の用を達し、自由自在、互いに人の妨げをなさずしておのおの安楽にこの世を渡らしめ給うの趣意なり。されども今、広くこの人間世界を見渡すに、かしこき人あり、おろかなる人あり、貧しきもあり、富めるもあり、貴人もあり、下人もありて、その有様雲と泥どろとの相違あるに似たるはなんぞや。その次第はなはだ明らかなり。『実語教じつごきょう』に、「人学ばざれば智なし、智なき者は愚人なり」とあり。されば賢人と愚人との別は学ぶと学ばざるとによりてできるものなり。また世の中にむずかしき仕事もあり、やすき仕事もあり。そのむずかしき仕事をする者を身分重き人と名づけ、やすき仕事をする者を身分軽き人という。すべて心を用い、心配する仕事はむずかしくして、手足を用うる力役りきえきはやすし。ゆえに医者、学者、政府の役人、または大なる商売をする町人、あまたの奉公人を召し使う大百姓などは、身分重くして貴き者と言うべし。
　身分重くして貴ければおのずからその家も富んで、下々しもじもの者より見れば及ぶべからざるようなれども、その本もとを尋ぬればただその人に学問の力あるとなきとによりてその相違もできたるのみにて、天より定めたる約束にあらず。諺ことわざにいわく、「天は富貴を人に与えずして、これをその人の働きに与うるものなり」と。されば前にも言えるとおり、人は生まれながらにして貴賤・貧富の別なし。ただ学問を勤めて物事をよく知る者は貴人となり富人となり、無学なる者は貧人となり下人げにんとなるなり。
　学問とは、ただむずかしき字を知り、解げし難き古文を読み、和歌を楽しみ、詩を作るなど、世上に実のなき文学を言うにあらず。これらの文学もおのずから人の心を悦よろこばしめずいぶん調法なるものなれども、古来、世間の儒者・和学者などの申すよう、さまであがめ貴とうとむべきものにあらず。古来、漢学者に世帯持ちの上手なる者も少なく、和歌をよくして商売に巧者なる町人もまれなり。これがため心ある町人・百姓は、その子の学問に出精するを見て、やがて身代を持ち崩すならんとて親心に心配する者あり。無理ならぬことなり。畢竟ひっきょうその学問の実に遠くして日用の間に合わぬ証拠なり。
　されば今、かかる実なき学問はまず次にし、もっぱら勤むべきは人間普通日用に近き実学なり。譬たとえば、いろは四十七文字を習い、手紙の文言もんごん、帳合いの仕方、算盤そろばんの稽古、天秤てんびんの取扱い等を心得、なおまた進んで学ぶべき箇条ははなはだ多し。地理学とは日本国中はもちろん世界万国の風土ふうど道案内なり。究理学とは天地万物の性質を見て、その働きを知る学問なり。歴史とは年代記のくわしきものにて万国古今の有様を詮索する書物なり。経済学とは一身一家の世帯より天下の世帯を説きたるものなり。修身学とは身の行ないを修め、人に交わり、この世を渡るべき天然の道理を述べたるものなり。
　これらの学問をするに、いずれも西洋の翻訳書を取り調べ、たいていのことは日本の仮名にて用を便じ、あるいは年少にして文才ある者へは横文字をも読ませ、一科一学も実事を押え、その事につきその物に従い、近く物事の道理を求めて今日の用を達すべきなり。右は人間普通の実学にて、人たる者は貴賤上下の区別なく、みなことごとくたしなむべき心得なれば、この心得ありて後に、士農工商おのおのその分を尽くし、銘々の家業を営み、身も独立し、家も独立し、天下国家も独立すべきなり。
　学問をするには分限を知ること肝要なり。人の天然生まれつきは、繋つながれず縛られず、一人前いちにんまえの男は男、一人前の女は女にて、自由自在なる者なれども、ただ自由自在とのみ唱えて分限ぶんげんを知らざればわがまま放蕩に陥ること多し。すなわちその分限とは、天の道理に基づき人の情に従い、他人の妨げをなさずしてわが一身の自由を達することなり。自由とわがままとの界さかいは、他人の妨げをなすとなさざるとの間にあり。譬たとえば自分の金銀を費やしてなすことなれば、たとい酒色に耽ふけり放蕩を尽くすも自由自在なるべきに似たれども、けっして然しからず、一人の放蕩は諸人の手本となり、ついに世間の風俗を乱りて人の教えに妨げをなすがゆえに、その費やすところの金銀はその人のものたりとも、その罪許すべからず。
　　このたび余輩の故郷中津に学校を開くにつき、学問の趣意を記して旧ふるく交わりたる同郷の友人へ示さんがため一冊を綴りしかば、或る人これを見ていわく、「この冊子をひとり中津の人へのみ示さんより、広く世間に布告せばその益もまた広かるべし」との勧めにより、すなわち慶応義塾の活字版をもってこれを摺すり、同志の一覧に供うるなり。
\end{breakbox}
\end{document}

